\section{这门课试图回答什么}

CS329A 关注的不是如何让一个语言模型在单次调用里回答得更像人，而是如何构造一个\textbf{能利用额外计算、外部工具和环境反馈持续提升表现的系统}。课程把过去几年的能力扩展路径串成一条连续主线：

\begin{enumerate}
    \item \textbf{预训练扩展（pre-training scaling）}：增加参数、数据与训练计算；
    \item \textbf{后训练（post-training）}：通过指令数据、偏好数据和强化学习改变模型行为；
    \item \textbf{测试时扩展（test-time scaling）}：在参数固定后，通过采样、搜索、验证和更长推理换取更高成功率；
    \item \textbf{Agent 系统}：把模型放入行动--反馈闭环，让它调用工具、读写环境、规划并纠错；
    \item \textbf{自我改进}：把生成、验证、反馈和训练重新连接成可迭代的闭环。
\end{enumerate}

\begin{importantbox}{本讲的总问题}
当单纯扩大预训练规模的边际收益开始下降时，能力提升还能从哪里来？本讲给出的答案是：\textbf{把计算从训练阶段扩展到推理阶段，并把单次模型调用升级为带验证与环境反馈的 Agent 工作流。}
\end{importantbox}

\subsection{学习目标}

完成本讲后，读者应能：

\begin{itemize}
    \item 解释参数、数据和训练计算为何形成预训练 scaling laws；
    \item 区分 zero-shot、few-shot、chain-of-thought 与 instruction tuning；
    \item 说明 RLHF 如何把人类偏好压缩为可优化的奖励信号；
    \item 推导 repeated sampling 的成功概率，并理解 generation--verification gap；
    \item 区分 LLM、agentic workflow 与自主 Agent；
    \item 用统一的“生成--验证--反馈--更新”视角理解整门课程。
\end{itemize}

\subsection{本章小结}

CS329A 的对象是完整系统：基础模型提供候选能力，推理计算扩大候选覆盖，验证器负责选择，工具与环境提供真实反馈，训练或记忆机制再把反馈沉淀为下一轮能力。

